\documentclass[11pt]{article}
\usepackage[utf8]{inputenc}
\usepackage{amsmath}
\usepackage{amssymb}
\usepackage{geometry}
\usepackage{booktabs}
\usepackage{hyperref}
\geometry{margin=1in}

\hypersetup{colorlinks=true, linkcolor=blue, urlcolor=blue}

\title{MethylClassifier Theoretical Foundation}
\author{MethylPipeline}
\date{}

\begin{document}
\maketitle
\tableofcontents
\newpage

% -----------------------------------------------------------------------
\section{Goal}

\texttt{MethylClassifier} assigns DNA methylation samples to biological
classes (e.g.\ healthy vs.\ cancer) using methylation at \emph{differentially
methylated positions} (DMPs).  It uses a \textbf{Bayesian Naive Bayes
framework with ECDF-based likelihoods}: no parametric distribution (Beta,
Normal, Beta-Binomial, or Beta-Mixture) is used.

Class-wise centroids provide an empirical distribution (ECDF) per position
via \texttt{binned\_stats} (\texttt{bin\_edges}, \texttt{bin\_counts}).  The
likelihood is the PDF from that ECDF, obtained by PCHIP spline interpolation.
DMP positions and centroid parameters come from \texttt{MethylDetector}
output.

% -----------------------------------------------------------------------
\section{Notation}

Let:
\begin{itemize}
  \item \(K\) --- number of classes (typically \(K=2\): control and disease).
  \item \(I\) --- number of selected DMP positions.
  \item \(x = (x_1,\ldots,x_I) \in [0,1]^I\) --- methylation fractions at
        DMP positions for a new sample.
  \item \(F_{k,i}(x)\) --- the PCHIP CDF at DMP position \(i\) for class
        \(k\), built from centroid \texttt{bin\_counts}.
  \item \(f_{k,i}(x) = F_{k,i}'(x)\) --- the corresponding PDF (spline
        derivative).
  \item \(w_i\) --- per-DMP weight (e.g.\ normalised \texttt{effect\_size}).
\end{itemize}

% -----------------------------------------------------------------------
\section{Per-Position Likelihood}

For DMP position \(i\) and class \(k\), the per-position likelihood is

\begin{equation}
  P(x_i \mid \text{class}\,k) = f_{k,i}(x_i),
  \label{eq:perpos_likelihood}
\end{equation}

where \(f_{k,i}\) is derived from the PCHIP spline fitted to the class-$k$
centroid histogram at position $i$.  A small floor is applied to avoid
\(\log(0)\): \(f_{k,i}(x) \geq 10^{-300}\).

% -----------------------------------------------------------------------
\section{Naive Bayes Likelihood over All Positions}

The Naive Bayes assumption asserts that methylation values at different DMP
positions are conditionally independent given the class:

\begin{equation}
  P(x \mid \text{class}\,k)
  = \prod_{i=1}^{I} P(x_i \mid \text{class}\,k)
  = \prod_{i=1}^{I} f_{k,i}(x_i).
\end{equation}

% -----------------------------------------------------------------------
\section{Weighted Log-Likelihood}

In practice the classifier uses \textbf{weighted log-likelihoods} for
numerical stability:

\begin{equation}
  \log L(\text{class}\,k \mid x)
  = \sum_{i=1}^{I} w_i \cdot \log f_{k,i}(x_i),
  \label{eq:loglik}
\end{equation}

where \(w_i = \mathrm{effect\_size}_i / \max_j(\mathrm{effect\_size}_j)\)
so that positions with larger biological effect contribute more.

% -----------------------------------------------------------------------
\section{Posterior Probability}

By Bayes' theorem with uniform priors:

\begin{equation}
  P(\text{class}\,k \mid x)
  \propto P(\text{class}\,k)\cdot P(x \mid \text{class}\,k)
  \propto \exp\!\bigl(\log L(\text{class}\,k \mid x)\bigr).
\end{equation}

The prior \(P(\text{class}\,k)\) is taken as uniform, so prediction depends
only on the log-likelihood term.

% -----------------------------------------------------------------------
\section{Temperature Scaling and Softmax}

A \emph{temperature} parameter \(T \geq 0.1\) controls the sharpness of the
posterior:

\begin{equation}
  P(\text{class}\,k \mid x)
  = \frac{\exp\!\left(\dfrac{\log L(\text{class}\,k \mid x)}{T}\right)}
         {\displaystyle\sum_{k'} \exp\!\left(\dfrac{\log L(\text{class}\,k' \mid x)}{T}\right)}.
  \label{eq:softmax}
\end{equation}

\begin{itemize}
  \item \(T \to 1^+\): posterior concentrates on the most likely class.
  \item \(T \gg 1\): posteriors become more uniform (softer decisions).
\end{itemize}

The default temperature is \(T = 2.0\).

% -----------------------------------------------------------------------
\section{Prediction}

The predicted class is

\begin{equation}
  \hat{k} = \arg\max_k P(\text{class}\,k \mid x).
\end{equation}

% -----------------------------------------------------------------------
\section{Optional Extensions}

\paragraph{Platt Calibration.}

Optionally, a logistic regression calibrator (Platt scaling) is fitted on
calibration data to improve the probability calibration of the raw posteriors
from Equation~\eqref{eq:softmax}.

\paragraph{Missing Positions.}

When a DMP position is not covered in a new sample (\(x_i\) is NaN), the
position is excluded from the sum in Equation~\eqref{eq:loglik}: its weight
contributes \(0\) to both classes.  This is equivalent to treating the
missing position as uninformative rather than imputing a neutral value.

% -----------------------------------------------------------------------
\section{Summary}

\begin{table}[h]
\centering
\begin{tabular}{@{}ll@{}}
\toprule
Aspect & Role \\
\midrule
Training & None in MethylClassifier; DMPs and centroids come from MethylDetector. \\
Likelihood & ECDF only: per-position PDF from centroid \texttt{bin\_counts}. \\
Combination & Weighted Naive Bayes log-sum over DMP positions. \\
Weights & Normalised \texttt{effect\_size} from MethylDetector. \\
Posterior & Temperature-scaled softmax over classes. \\
Output & Posterior probabilities per class; prediction = argmax. \\
\bottomrule
\end{tabular}
\caption{Summary of the MethylClassifier framework.}
\end{table}

\end{document}
